\documentclass[11pt, a4paper]{report}

\usepackage[utf8]{inputenc}
\usepackage[T1]{fontenc}
\usepackage[french]{babel}

% Appel des packages dans le dossier preambule
\def\packagepath{../../../../preambule}   % path du package princial

\usepackage{\packagepath/page_layout} % <--- Nouveau
\usepackage{\packagepath/config_info}
\usepackage{\packagepath/cadres_cours}
\usepackage{\packagepath/zones_reponse}
\usepackage{\packagepath/config_ds}
\usepackage{\packagepath/config_maths}

% --- PILOTAGE ---
\def\visibleornot{invisible}    % visible or invisible


\def\documentpath{./Sources_Latex}
\graphicspath{{./Images/}}


\def\ptsexoA{0}

\begin{document}

%%%%%%%%%%%%%%%%%%%%%%%%%%%
\def\level{BTS}  % Classe à droite
\def\course{CIEL 1}
\def\chaptername{Intégration} % Titre au centre
\def\docname{C107~--~TP1}        % Identifiant à gauche

\pagestyle{DS_FP}
%%%%%%%%%%%%%%%%%%%%%%%%%%%%%%%%%

\NomPrenomNote{}

%%=============================================================

\section*{PARTIE A\@: Mise en place de la fonction à étudier}
\medskip

L'objectif de cette partie est l'étude de la fonction $f$ définie sur $\R$ par:

\[f(x)=-0,5x^2+2x+1\]


\begin{enumerate}
    \item À l'aide de Géogebra, tracer la courbe représentative de la fonction $f$. 
    \item Créer deux curseurs $a$ et $b$ entre 0 et 5 avec un pas de 0,1 et donner les valeurs suivantes: $a=1$ et $b=4$. 
\end{enumerate}

\section*{Partie B\@: Somme de Riemann (Approximation)}

On souhaite diviser l'intervalle $\ff{a}{b}$ en $n$ sous-intervalles de même longueur. 

\begin{enumerate}
    \item Créer un curseur $n$ compris entre 1 et 100 avec un pas de 1. 
    \item Saisissez la commande \verb|S_inf = SommeInférieure(f, a, b, n)|. Changer sa couleur.
    \item Saisissez la commande \verb|S_sup = SommeSupérieure(f, a, b, n)|. Changer sa couleur.
    \item Observer les valeurs de $S_{\inf}$ et $S_{\sup}$ pour différentes valeurs de $n$.  
    \item Que peut-on conjecturer sur la valeur de l'aire sous la courbe de $f$ entre $a$ et $b$ lorsque $n$ tend vers l'infini? 
    \begin{blocvisible}[height=3]
                Lorsque $n$ tend vers l'infini, les sommes de Riemann $S_{\inf}$ et $S_{\sup}$ convergent vers la même valeur, qui correspond à l'aire sous la courbe de la fonction $f$ entre les points $a$ et $b$. On peut conjecturer que cette aire est égale à la limite de ces sommes lorsque $n$ tend vers l'infini.
                \[\lim_{n \to \infty} S_{\inf} = \lim_{n \to +\infty} S_{\sup} = \lim_{n \to \infty} \int_{a}^{b} f(x) \dx= \text{Aire sous la courbe de } f \text{ entre } a \text{ et } b\]        
    \end{blocvisible}     
\end{enumerate}

\section*{Partie C\@: Calcul de l'aire sous la courbe à l'aide de primitives}

\begin{enumerate}
    \item Déterminer une primitive $F$ de la fonction $f$. 
        \begin{blocvisible}[height=2]
            \[F(x)=-0,5 \cdot \frac{x^3}{3} + 2 \cdot \frac{x^2}{2} + 1 \cdot x + C= -\frac{x^3}{6} + x^2 + x + C, \quad C \in \R\]
        \end{blocvisible}

    \item Calculer $F(b)-F(a)$ et interpréter ce résultat. 
        \begin{blocvisible}[height=3]
                En calculant $F(b)-F(a)$, on obtient la valeur de l'aire sous la courbe de la fonction $f$ entre les points $a$ et $b$. 
                \[F(b)-F(a) = \int_{a}^{b} f(x) \dx=\int_{1}^{4} f(x) \dx= F(4)-F(1)=\frac{28}{3} - \frac{11}{6} = \frac{56 - 11}{6} = \frac{45}{6} = \frac{15}{2}=7,5 \qquad \texttt{u.a}\]
        \end{blocvisible}

    \item Dans Géogebra, saisir la commande \verb|Aire = Intégrale(f, a, b)|. Comparer la valeur de $F(b)-F(a)$ avec celle de l'aire calculée dans Géogebra. 
        \begin{blocvisible}[height=2]
                La valeur de $F(b)-F(a)$ est égale à l'aire calculée dans Géogebra, ce qui confirme que le calcul de l'aire sous la courbe de la fonction $f$ entre les points $a$ et $b$ est correct. 
        \end{blocvisible}

\pagebreak
\thispagestyle{DS_LP}

    \item Calculer la valeur moyenne de la fonction $f$ sur l'intervalle $\ff{a}{b}$ et interpréter ce résultat. 
        \begin{blocvisible}[height=4]
            La valeur moyenne de la fonction $f$ sur l'intervalle $\ff{a}{b}$ est donnée par la formule:
            \[\mu_{\ff{a}{b}}=\frac{1}{b-a} \int_{a}^{b} f(x) \dx = \frac{1}{4-1} \cdot \frac{15}{2} = \frac{1}{3} \cdot \frac{15}{2} = \frac{5}{2}=2,5\]
            Cette valeur correspond à la hauteur du rectangle dont l'aire est égale à l'aire sous la courbe de $f$ entre $a$ et $b$.
        \end{blocvisible}
    
    \item Créer avec Géogébra le rectangle $ABCD$ tel que $AB$ soit égal à la valeur moyenne de $f$ sur $\ff{a}{b}$ et que $AD$ soit égal à $b-a$. Vérifier que l'aire de ce rectangle est égale à l'aire sous la courbe de $f$ entre $a$ et $b$. 
        \begin{blocvisible}[height=4]
            En créant le rectangle $ABCD$ avec $AB$ égal à la valeur moyenne de $f$ sur $\ff{a}{b}$ (soit 2,5) et $AD$ égal à $b-a$ (soit 3), on obtient un rectangle dont l'aire est donnée par:
            \[\text{Aire}_{ABCD} = AB \cdot AD = 2,5 \cdot 3 = 7,5\]
            Cette aire est égale à l'aire sous la courbe de $f$ entre les points $a$ et $b$, confirmant ainsi que le rectangle $ABCD$ a la même aire que celle sous la courbe de la fonction $f$ sur l'intervalle $\ff{a}{b}$.
        \end{blocvisible}
\end{enumerate}
\end{document}







